\documentclass[12pt,a4paper]{article}
\usepackage[margin=1in]{geometry}
\usepackage{enumitem}
\usepackage{titlesec}
\usepackage{parskip}
\usepackage{booktabs}
\usepackage{array}
\usepackage{amsmath}
\usepackage{fancyhdr}
\usepackage{hyperref}
\usepackage{xcolor}
\usepackage{listings}

% Header / Footer
\pagestyle{fancy}
\fancyhf{}
\rhead{PRMV Lab}
\lhead{Pattern Recognition \& Machine Vision}
\cfoot{\thepage}

% Section formatting
\titleformat{\section}{\large\bfseries}{Experiment \thesection:}{0.5em}{}[\titlerule]
\titleformat{\subsection}{\normalsize\bfseries}{}{0em}{}

% Code style
\lstset{
  basicstyle=\ttfamily\small,
  backgroundcolor=\color{gray!10},
  frame=single,
  breaklines=true,
  keywordstyle=\color{blue},
  commentstyle=\color{green!60!black},
  numbers=none
}

\hypersetup{colorlinks=true, linkcolor=blue, urlcolor=blue}

\title{\textbf{Pattern Recognition \& Machine Vision}\\[0.5em]
\large Laboratory Experiments Write-Up}
\author{PRMV Lab}
\date{}

\begin{document}

\maketitle
\tableofcontents
\newpage

%=============================================================
\section{Feature Extraction from Images and Signals}
%=============================================================

\subsection*{Aim}
To implement basic feature extraction techniques (mean, variance, and histogram features)
and visualize feature distributions using image or signal datasets.

\subsection*{Algorithm}
\begin{enumerate}[leftmargin=*, label=\arabic*.]
  \item Load all images from the \texttt{images/} directory using \texttt{glob}.
  \item Convert each image to grayscale using \texttt{cv2.IMREAD\_GRAYSCALE}.
  \item For each image, compute:
    \begin{itemize}
      \item \textbf{Mean}: \texttt{np.mean(img)} — average pixel intensity.
      \item \textbf{Variance}: \texttt{np.var(img)} — spread of pixel values.
      \item \textbf{Histogram}: \texttt{cv2.calcHist()} — 256-bin pixel intensity distribution.
    \end{itemize}
  \item Store features in separate lists.
  \item Plot bar charts for mean and variance per image, and a line plot for the average histogram.
\end{enumerate}

\subsection*{Output Obtained}
\begin{itemize}
  \item Bar chart of \textbf{mean pixel intensity} per image.
  \item Bar chart of \textbf{variance} per image.
  \item Line plot of the \textbf{average histogram} across all images (pixel intensity 0--255 on x-axis).
\end{itemize}

\noindent Sample feature values (approximate):
\begin{center}
\begin{tabular}{lcc}
\toprule
\textbf{Image} & \textbf{Mean (approx.)} & \textbf{Variance (approx.)} \\
\midrule
car.jpg       & 140 & 3500 \\
cat.jpg       & 160 & 3000 \\
dandelion.jpg & 170 & 2800 \\
dog.jpg       & 130 & 4000 \\
fish.jpg      & 150 & 3200 \\
\bottomrule
\end{tabular}
\end{center}

\subsection*{Result}
Feature extraction was successfully implemented on a set of grayscale images. Mean pixel intensity
reflected image brightness, variance represented contrast and texture richness, and the histogram
provided a complete distribution of pixel intensities. These features form the foundation for image
classification in subsequent experiments.

\newpage

%=============================================================
\section{Two-Class Gaussian Bayes Classifier}
%=============================================================

\subsection*{Aim}
To implement a Bayes classifier for two classes (Flower and Dog) using Gaussian distributions
to compute posterior probabilities and evaluate classification accuracy.

\subsection*{Algorithm}
\begin{enumerate}[leftmargin=*, label=\arabic*.]
  \item Load images from \texttt{flower/} (Class 0) and \texttt{dog/} (Class 1) directories.
  \item Convert each image to grayscale; extract two features: pixel mean and pixel standard deviation.
  \item Split the dataset: 80\% training, 20\% testing (\texttt{np.random.seed(42)}).
  \item For each class, estimate:
    \begin{itemize}
      \item Class mean vector $\boldsymbol{\mu}_c$ and covariance matrix $\boldsymbol{\Sigma}_c$.
      \item Prior probability $P(\omega_c)$ from class frequency.
    \end{itemize}
  \item Compute the multivariate Gaussian likelihood:
    \[
      p(\mathbf{x} \mid \omega_c) = \frac{1}{\sqrt{(2\pi)^D |\boldsymbol{\Sigma}_c|}}
      \exp\!\left(-\tfrac{1}{2}(\mathbf{x}-\boldsymbol{\mu}_c)^T \boldsymbol{\Sigma}_c^{-1}(\mathbf{x}-\boldsymbol{\mu}_c)\right)
    \]
  \item Compute posterior $\propto$ likelihood $\times$ prior; predict class as $\arg\max$.
  \item Evaluate: Accuracy, Precision, Recall, F1-Score, Confusion Matrix.
\end{enumerate}

\subsection*{Output Obtained}
\begin{itemize}
  \item Total images: 1545 \quad (Flower = 843, Dog = 702)
  \item Train: 1236 \quad Test: 309
\end{itemize}

\begin{center}
\begin{tabular}{lc}
\toprule
\textbf{Metric} & \textbf{Value} \\
\midrule
Test Accuracy  & 52.10\% \\
Precision      & 0.5140  \\
Recall         & 0.3642  \\
F1-Score       & 0.4264  \\
\bottomrule
\end{tabular}
\end{center}

\noindent Confusion Matrix:
\[
\begin{bmatrix} TN = 106 & FP = 52 \\ FN = 96 & TP = 55 \end{bmatrix}
\]

\noindent \textbf{Plots generated:} 1D class-conditional Gaussian PDFs (pixel mean and std dev),
and a 2D scatter plot of test samples with class mean markers.

\subsection*{Result}
A two-class Gaussian Bayes Classifier was implemented from scratch using pixel mean and standard
deviation as features. The classifier achieved an accuracy of 52.10\% on the Flower vs.\ Dog dataset.
The moderate accuracy is due to the simplicity of the features; the Bayesian framework with
multivariate Gaussian class-conditionals and estimated priors was successfully demonstrated.

\newpage

%=============================================================
\section{Linear Discriminant Classifier vs.\ Gaussian Bayes}
%=============================================================

\subsection*{Aim}
To implement a Linear Discriminant Classifier and compare it with the Bayesian classifier by
visualizing decision boundaries.

\subsection*{Algorithm}
\begin{enumerate}[leftmargin=*, label=\arabic*.]
  \item Load images from \texttt{fruit/} (Class 0) and \texttt{dog/} (Class 1); extract pixel mean and std dev.
  \item Split dataset: 80\% train, 20\% test.
  \item \textbf{LDA fit:}
    \begin{itemize}
      \item Compute class means $\boldsymbol{\mu}_c$ and priors $P(\omega_c)$.
      \item Compute pooled scalar variance: $\sigma^2 = \frac{\sum_c \sum_{\mathbf{x} \in \omega_c} \|\mathbf{x}-\boldsymbol{\mu}_c\|^2}{N \cdot D}$
      \item Compute weight vectors $\mathbf{w}_c = \boldsymbol{\mu}_c / \sigma^2$ and biases $w_{c0} = -\frac{\|\boldsymbol{\mu}_c\|^2}{2\sigma^2} + \ln P(\omega_c)$.
    \end{itemize}
  \item \textbf{LDA predict:} Classify as class 1 if $g_1(\mathbf{x}) > g_0(\mathbf{x})$, where $g_c(\mathbf{x}) = \mathbf{w}_c^T\mathbf{x} + w_{c0}$.
  \item \textbf{Gaussian Bayes:} Fit and predict as in Experiment 2.
  \item Evaluate both classifiers; plot side-by-side decision boundary maps.
\end{enumerate}

\subsection*{Output Obtained}

\begin{center}
\begin{tabular}{lcc}
\toprule
\textbf{Classifier} & \textbf{Accuracy} & \textbf{Boundary Type} \\
\midrule
Linear Discriminant (LDA) & 82.40\% & Straight line (linear) \\
Gaussian Bayes            & 94.72\% & Curved (quadratic)     \\
\bottomrule
\end{tabular}
\end{center}

\noindent LDA Parameters:
\begin{itemize}
  \item Shared variance $\sigma^2 = 539.8042$
  \item Boundary direction vector $(\mathbf{w}_1 - \mathbf{w}_0) = [-0.084,\ -0.028]$
  \item Boundary offset $b = 13.1076$
\end{itemize}

\noindent \textbf{Plot generated:} Side-by-side decision boundary maps showing LDA's straight boundary
vs.\ Bayes' curved boundary, overlaid with training and test data points.

\subsection*{Result}
Both classifiers were implemented and compared on the Fruit vs.\ Dog dataset. The Gaussian Bayes
classifier outperformed LDA (94.72\% vs.\ 82.40\%), because the two classes have significantly
different covariance structures. LDA's assumption of a shared isotropic variance is too restrictive,
while Bayes with class-specific covariances produces a more flexible decision boundary.

\newpage

%=============================================================
\section{Maximum Likelihood Estimation (MLE)}
%=============================================================

\subsection*{Aim}
To develop and evaluate parametric estimation techniques employing Maximum Likelihood
Estimation (MLE).

\subsection*{Algorithm}
\begin{enumerate}[leftmargin=*, label=\arabic*.]
  \item Generate $N = 200$ samples from $\mathcal{N}(\mu=50,\ \sigma=15)$ using \texttt{np.random.normal}.
  \item \textbf{Case 1 --- $\sigma$ known, estimate $\mu$:}
    \begin{itemize}
      \item MLE estimator: $\hat{\mu} = \frac{1}{N}\sum_{k=1}^{N} x_k$
      \item Compute log-likelihood $l(\hat{\mu})$ and verify $\frac{\partial l}{\partial \mu} = 0$.
    \end{itemize}
  \item \textbf{Case 2 --- Both $\mu$ and $\sigma^2$ unknown:}
    \begin{itemize}
      \item MLE: $\hat{\mu} = \frac{1}{N}\sum x_k$, \quad $\hat{\sigma}^2 = \frac{1}{N}\sum(x_k - \hat{\mu})^2$
      \item Compute log-likelihood $l(\hat{\mu}, \hat{\sigma}^2)$.
    \end{itemize}
  \item Plot true Gaussian PDF vs.\ MLE-estimated PDF for both cases side by side.
\end{enumerate}

\subsection*{Output Obtained}

\begin{center}
\begin{tabular}{lccc}
\toprule
\textbf{Case} & \textbf{Parameter} & \textbf{True Value} & \textbf{MLE Estimate} \\
\midrule
\multirow{1}{*}{Case 1} & $\mu$ & 50.0 & \textbf{49.3884} \\
\midrule
\multirow{2}{*}{Case 2} & $\mu$    & 50.0 & \textbf{49.3884} \\
                         & $\sigma$ & 15.0 & \textbf{13.9301} \\
\bottomrule
\end{tabular}
\end{center}

\begin{itemize}
  \item Log-likelihood (Case 1): $l(\hat{\mu}) = -811.6412$
  \item Log-likelihood (Case 2): $l(\hat{\mu}, \hat{\sigma}) = -810.5981$
\end{itemize}

\noindent \textbf{Plot generated:} \texttt{mle\_results.png} --- two panels comparing the true
Gaussian PDF against the MLE-estimated PDFs for Case 1 and Case 2.

\subsection*{Result}
MLE was successfully applied to estimate the parameters of a Gaussian distribution from 200
synthetic samples. The estimated $\hat{\mu} = 49.39$ closely matches the true $\mu = 50$.
The slight underestimation of $\sigma$ (13.93 vs.\ 15.0) is due to the biased MLE variance
estimator, which divides by $N$ rather than $N-1$.

\newpage

%=============================================================
\section{Non-Parametric Density Estimation using k-NN}
%=============================================================

\subsection*{Aim}
To design and analyze non-parametric density estimation methods based on
k-Nearest Neighbours (k-NN).

\subsection*{Algorithm}
\begin{enumerate}[leftmargin=*, label=\arabic*.]
  \item Generate $N = 300$ samples from a Gaussian mixture:
    $0.4 \cdot \mathcal{N}(20,5^2) + 0.6 \cdot \mathcal{N}(50,8^2)$.
  \item Define a fine evaluation grid over the data range (500 points).
  \item For each query point $x$ and each $k \in \{1, 5, 15, 30\}$:
    \begin{itemize}
      \item Compute distances: $|x - x_i|$ for all $x_i$ in data.
      \item Find $d_k$ = distance to the $k$-th nearest neighbour.
      \item Estimate density: $\hat{p}(x) = \dfrac{k}{n \cdot 2d_k}$
    \end{itemize}
  \item Compute MSE between estimated density and the true mixture PDF.
  \item Plot estimated density vs.\ true PDF for each value of $k$.
\end{enumerate}

\subsection*{Output Obtained}

\begin{center}
\begin{tabular}{ccc}
\toprule
\textbf{k} & \textbf{MSE} & \textbf{Estimate Quality} \\
\midrule
1  & 0.047246 & Very spiky / noisy \\
5  & 0.000154 & Noisy but improving \\
15 & 0.000026 & Smooth \\
30 & 0.000015 & Smoothest, best fit \\
\bottomrule
\end{tabular}
\end{center}

\noindent \textbf{Plot generated:} \texttt{knn\_results.png} --- 4 panels showing the data
histogram, true PDF, and k-NN estimate for each value of $k$.

\subsection*{Result}
Non-parametric k-NN density estimation was successfully implemented for a 1D Gaussian mixture.
As $k$ increases, the density estimate becomes progressively smoother and converges closer
to the true PDF (lower MSE). Very small $k$ (e.g., $k=1$) produces spiky, discontinuous
estimates, while larger $k$ (e.g., $k=30$) gives a smooth, accurate approximation of the
underlying distribution.

\newpage

%=============================================================
\section{Dimensionality Reduction using PCA}
%=============================================================

\subsection*{Aim}
To perform dimensionality reduction using Principal Component Analysis (PCA).

\subsection*{Algorithm}
\begin{enumerate}[leftmargin=*, label=\arabic*.]
  \item Load the Iris dataset (150 samples, 4 features, 3 classes) using \texttt{sklearn}.
  \item \textbf{Center the data:} $\mathbf{X}_{std} = \mathbf{X} - \bar{\mathbf{X}}$
  \item \textbf{Compute covariance matrix:} $\mathbf{C} = \text{np.cov}(\mathbf{X}_{std},\ \text{rowvar=False})$ \quad (4$\times$4)
  \item \textbf{Eigendecomposition:} \texttt{np.linalg.eigh(C)} --- obtain eigenvalues and eigenvectors.
  \item Sort eigenvectors in descending order of eigenvalue.
  \item Compute variance explained per component: $\frac{\lambda_i}{\sum \lambda} \times 100\%$.
  \item Select top-2 eigenvectors to form projection matrix $\mathbf{W}$ (4$\times$2).
  \item Project data to 2D: $\mathbf{X}_{2D} = \mathbf{X}_{std} \cdot \mathbf{W}$
  \item Plot: eigenvalue bar chart, cumulative variance curve, 2D scatter of projected data.
\end{enumerate}

\subsection*{Output Obtained}

\begin{center}
\begin{tabular}{lccr}
\toprule
\textbf{Component} & \textbf{Eigenvalue} & \textbf{Variance Explained} & \textbf{Cumulative} \\
\midrule
PC1 & 4.2282 & 92.46\% & 92.46\% \\
PC2 & 0.2427 &  5.31\% & 97.77\% \\
PC3 & 0.0782 &  1.71\% & 99.48\% \\
PC4 & 0.0238 &  0.52\% & 100.00\% \\
\bottomrule
\end{tabular}
\end{center}

\begin{itemize}
  \item \textbf{Variance retained with PC1 + PC2: 97.77\%}
  \item PC1 eigenvector: $[-0.3614,\ 0.0845,\ -0.8567,\ -0.3583]$
  \item PC2 eigenvector: $[\ 0.6566,\ 0.7302,\ -0.1734,\ -0.0755]$
  \item Reduced data sample (first row): $[2.6841,\ 0.3194]$
\end{itemize}

\noindent \textbf{Plot generated:} \texttt{pca\_results.png} --- eigenvalue bar chart, cumulative
variance curve with 95\% threshold line, and 2D scatter plot of three Iris classes.

\subsection*{Result}
PCA was successfully applied to reduce the 4-dimensional Iris dataset to 2 dimensions while
retaining 97.77\% of the total variance. The 2D scatter plot shows that Setosa is clearly
separated from Versicolor and Virginica. PC1 alone captures 92.46\% of the total variation,
confirming that the data has a dominant principal direction corresponding to petal dimensions.

\newpage

%=============================================================
\section{Object Detection using YOLOv8}
%=============================================================

\subsection*{Aim}
To develop and demonstrate an object detection system.

\subsection*{Algorithm}
\begin{enumerate}[leftmargin=*, label=\arabic*.]
  \item Load a test image using \texttt{skimage.data.cat()} (built-in RGB cat image).
  \item Load the pre-trained YOLOv8s model: \texttt{YOLO('yolov8s.pt')}.
  \item Run inference with thresholds:
    \begin{itemize}
      \item \textbf{conf = 0.35} --- minimum confidence score to retain a detection.
      \item \textbf{iou = 0.50} --- IoU threshold for Non-Maximum Suppression (NMS).
      \item \textbf{imgsz = 640} --- input image size for the model.
    \end{itemize}
  \item For each detected bounding box, extract:
    \begin{itemize}
      \item Coordinates $[x_1, y_1, x_2, y_2]$ via \texttt{box.xyxy[0]}.
      \item Confidence score via \texttt{box.conf[0]}.
      \item Class label via \texttt{result.names[int(box.cls[0])]}.
    \end{itemize}
  \item Draw annotated bounding boxes using \texttt{result.plot()}.
  \item Display the annotated image using \texttt{matplotlib}.
\end{enumerate}

\subsection*{Output Obtained}

\begin{verbatim}
Successfully detected 1 object(s).

Object 1: Cat | Confidence: 0.87
         | Box: (x1=22.3, y1=5.8, x2=595.0, y2=478.0)
\end{verbatim}

\begin{center}
\begin{tabular}{ll}
\toprule
\textbf{Parameter} & \textbf{Value} \\
\midrule
Model              & YOLOv8s (pre-trained on COCO) \\
Test image         & \texttt{skimage.data.cat()} \\
Confidence threshold & 0.35 \\
IoU threshold      & 0.50 \\
Detected class     & Cat \\
Confidence score   & $\approx 0.87$ \\
Objects detected   & 1 \\
\bottomrule
\end{tabular}
\end{center}

\noindent \textbf{Visual output:} Image of a cat with a bounding box drawn around it,
labelled ``cat 0.87''.

\subsection*{Result}
An object detection system was successfully implemented using the pre-trained YOLOv8s model.
The model correctly detected the cat in the test image with high confidence ($\approx 0.87$).
The bounding box accurately localised the object within the image. YOLO's single-pass
architecture enables real-time object detection, making it highly suitable for practical
object recognition applications.

\newpage

%=============================================================
\section{Object Recognition and Motion Tracking using Video}
%=============================================================

\subsection*{Aim}
To design an object recognition and motion tracking system using video sequences.

\subsection*{Algorithm}
\begin{enumerate}[leftmargin=*, label=\arabic*.]
  \item Open the webcam using \texttt{cv2.VideoCapture(0)}.
  \item Define HSV colour range for the target colour (orange):
    \begin{itemize}
      \item Lower bound: \texttt{[5, 120, 80]}
      \item Upper bound: \texttt{[25, 255, 255]}
    \end{itemize}
  \item For each frame captured from the webcam:
    \begin{enumerate}[label=\alph*.]
      \item Convert frame from BGR to HSV: \texttt{cv2.cvtColor(frame, COLOR\_BGR2HSV)}.
      \item Create binary mask of the target colour: \texttt{cv2.inRange(hsv, lower, upper)}.
      \item Find contours in the mask: \texttt{cv2.findContours(...)}.
      \item Select the largest contour (area $> 500$ pixels) as the tracked object.
      \item Compute bounding rectangle: \texttt{cv2.boundingRect(largest)}.
      \item Draw the bounding box and label on the original frame.
      \item Display the annotated frame using \texttt{cv2.imshow()}.
    \end{enumerate}
  \item Repeat for every frame until the user presses \texttt{q} to quit.
  \item Release camera and close windows.
\end{enumerate}

\subsection*{Output Obtained}

\begin{center}
\begin{tabular}{ll}
\toprule
\textbf{Parameter} & \textbf{Value} \\
\midrule
Input source          & Webcam (\texttt{VideoCapture(0)}) \\
Target colour         & Orange \\
HSV lower bound       & \texttt{[5, 120, 80]} \\
HSV upper bound       & \texttt{[25, 255, 255]} \\
Minimum contour area  & 500 pixels$^2$ \\
Bounding box colour   & Orange (\texttt{BGR: 0, 140, 255}) \\
Tracking output       & Live bounding box + ``orange'' label \\
\bottomrule
\end{tabular}
\end{center}

\noindent \textbf{Visual output:} A live camera window titled
\textit{``Orange Tracker [q to quit]''}. When an orange object is held in front of the
camera, a bounding rectangle appears around it in real time with the label ``orange''
displayed above the box. Pressing \texttt{q} closes the window.

\subsection*{Result}
A real-time colour-based object tracking system was successfully implemented using OpenCV.
By converting each video frame to HSV colour space and applying colour thresholding followed
by contour detection, the system accurately locates and tracks an orange object in a live
webcam feed. The bounding box updates correctly in every frame, demonstrating a practical
object recognition and motion tracking pipeline.

\end{document}
